Understanding the current capabilities in AI Scientist systems  (autoresearch), autonomous agents capable of conducting scientific research, is crucial for promoting sustainable, AI-driven scientific progress.
Nevertheless, developers must remain conscious of the potential risks these systems pose to the academic ecosystem and commit to advancing them responsibly.
Since 2025, a new venue dedicated to evaluating AI-driven scientific contributions, the Agents4Science conference~\citep{agents4science, bianchi2025exploring}, has emerged. 
Through such a platform, developers of AI Scientist systems are encouraged to engage in responsible research and development, ensuring both the protection of the academic ecosystem and the long-term sustainability of scientific progress.

In recent years, several works have explored the concept of AI Scientists~\citep{lu2024ai, weng2024cycleresearcher, yamada2025ai, tang2025ai}. However, the quality of research papers produced by these systems remains insufficient. One major reason is that the problem setting of achieving fully automated science is overly ambitious and often lacks clearly defined scientific goals for AI Scientists. Without a specific goal, these systems tend to generate undirected discoveries that appear to lack genuine scientific value. Another limitation is that current systems are limited to small-scale code experiments~\citep{zhu2025ai, lu2024ai, yamada2025ai}, lacking the scale and complexity needed for meaningful science. Achieving real scientific contributions requires not just ideas but strong implementation capability to handle complex codebases.

As an initial step toward enabling AI Scientists to produce genuine scientific value, we can take inspiration from how student researchers begin their research. When a student first joins a research lab, a common and meaningful process is as follows: the mentor assigns a key paper, the student analyzes its limitations, proposes an improvement hypothesis, implements the idea on the baseline code, iteratively experiments until improvements are achieved, and finally writes a paper summarizing the results. Through this process, the student learns the fundamental workflow of scientific research and gains the skills and experience needed for more creative work later on. Also, improving a baseline method is not only an important stage in early research training but also a valuable research goal in many fields where advancing task performance remains a central scientific pursuit.



\begin{figure*}[t]
    \centering
    \resizebox{0.95\textwidth}{!}{%
        \includegraphics{images/teaser.pdf}
    }
    \caption{\textbf{Jr. AI Scientist Workflow.} We provide the baseline paper, its LaTeX source files, and the associated codebase. By effectively utilizing these resources across all phases, the system significantly improves the quality of the generated paper.}
    \label{Fig:teaser}
\vspace{-1em}
\end{figure*}

In this paper, we introduce Jr. AI Scientist, a new AI Scientist that mimics the essential research workflow of a novice student researcher: Starting from a baseline paper, it identifies key limitations, proposes an improvement hypothesis, iteratively conduct experiments until improvements are achieved, and writes a paper with the results.
Specifically, A Jr. AI Scientist is defined as an AI Scientist that is given baseline resources and focuses on extending the baseline. This setting is novice student analog, where a student builds upon an existing paper which the mentor gave.
The workflow of Jr. AI Scientist is shown in \Cref{Fig:teaser}. We provide the baseline paper, its LaTeX source files, and the associated codebase for each stage.
\Cref{tab:ai_scientist_problem_setting} shows a comparison of the problem setting with existing research.
This problem setting reframes previously ambitious goals into a more specific objective, providing a clear optimization direction for the AI Scientist. Moreover, because the framework operates on the actual codebase of baseline papers, it can generate results with genuine scientific value. These aspects collectively represent an essential first step toward the autonomous generation of reliable and high-quality research papers.


Our Jr. AI Scientist consists of three main components: (1) automatic idea generation based on the limitations of a given paper, (2) automatic implementation and iteratively experiments until improvements are achieved, and (3) automatic writing of a research paper based on the obtained results.
This system is built upon AI Scientist v2~\citep{yamada2025ai}, but our work differs from prior studies~\citep{lu2024ai, weng2024cycleresearcher, yamada2025ai, tang2025ai} in several key aspects: First, by leveraging the latest coding agents (\eg~Claude Code~\citep{claude_code}), our system can perform meaningful improvements and edits on real multi-file codebases, which were challenging for previous AI Scientist systems. Second, by incorporating the full set of resources from a given baseline paper, the system exploits all available artifacts such as LaTeX sources, PDFs, and codebases, thereby substantially improving the scope and quality of every stage in the research pipeline. Finally, by refining every stage of the process, our framework enables the autonomous generation of research papers that are both higher in quality and more trustworthy.

\begin{table}[t]
\centering
\small
\caption{
Comparison of the starting point, code complexity, and review scores among existing AI Scientist systems.
Previous methods often made overly ambitious assumptions in their problem formulation and were limited to handling only simple, single-file codebases, which resulted in significantly lower review scores.
In contrast, our system can substantially improve review scores by utilizing the baseline paper and its associated codebase.
}
\label{tab:ai_scientist_problem_setting}
\vspace{5pt}
\begin{adjustbox}{width=\textwidth}
\begin{tabular}{lcccccc}
\toprule
\textbf{AI Scientist Systems} & \textbf{Starting Point} & \textbf{Code Complexity} & \textbf{Review Score} \\
\midrule
AI Scientist-v1~\citep{lu2024ai} & Template code  & Single file & 3.30 \\
AI Scientist-v2~\citep{yamada2025ai} & General idea  &  Single file & 2.75 \\
AI Researcher~\citep{tang2025ai} & 15-20 existing papers &  Multiple files  &  3.25 \\
Jr. AI Scientist (ours) & One baseline paper and code  & Multiple files & \textbf{5.75} \\
% DeepScientist (Ours) & 5  & \textbf{2.90} & \textbf{2.90} & \textbf{2.90} & \textbf{5.90}&\textbf{60\%} \\
\bottomrule
\end{tabular}
\end{adjustbox}
\label{table:max_score}
\end{table}

For the baseline papers, we selected papers for which we obtained permission from the original authors. Specifically, we used three papers: NeurIPS 2023 paper~\citep{miyai2023locoop} and IJCV 2025 paper~\citep{miyai2025gl} on out-of-distribution (OOD) detection (a task that aims to detect semantic classes outside the predefined set of semantic classes)~\citep{msp17iclr, yang2024generalized}, and an ICLR 2025 spotlight paper~\citep{zhang2024min} on pre-training data detection for large language models (LLMs). Refer to \Cref{subsec:baseline_selection} for the detailed rationale behind the selection of these papers. 
For the evaluation, we conducted three evaluations: (1) an automated assessment using DeepReviewer~\citep{zhu2025deepreview}, (2) an author-led evaluation, and (3) submission to the Agents4Science conference~\citep{agents4science}.
DeepReviewer automatically compared our generated papers with existing AI-generated works to assess overall quality.
The author evaluation examined the outputs for hallucinations or fabricated content.
Finally, the Agents4Science~\citep{agents4science} platform provides rigorous evaluation and feedback from the community platform.

Through our experiments, the Jr. AI Scientist successfully generated new research papers that build upon the above top-tier venues works by proposing and implementing novel algorithms.
As for the evaluation using DeepReviewer~\citep{zhu2025deepreview}, the papers generated by Jr. AI Scientist achieved substantially higher review scores compared to the existing AI-generated papers. Therefore, our Jr. AI Scientist can be regarded as the most capable autonomous AI Scientist.
However, we also observed that Jr. AI Scientist still exhibits some failures and unresolved challenges through the author evaluation and Agents4Science conference. To share these challenges and lessons with the research community, we analyze the feedback and evaluation results from the Agents4Science conference and include the author evaluation, which helps clarify what is required to further improve the quality of Jr. AI Scientist systems. 

Finally, we perform an in-depth report of the risks encountered during the development of our system. Although few existing studies have provided a comprehensive discussion of these issues, we believe that accurately documenting such risks is essential to avoid overestimating current AI Scientists' capabilities and to build a clear understanding of their remaining challenges. 
Our risk report highlights several critical issues, including the potential for review-score hacking and difficulties in ensuring proper citation,  interpreting results, and detecting fabricated descriptions.
We believe that these findings provide valuable guidance on the potential risks that exist both in the current AI Scientist research and as this field continues to grow. Through this comprehensive report, we aim to foster a deeper understanding of current AI Scientist systems and contribute to their safe and trustworthy development.

Our contributions are summarized as follows:
\begin{itemize}
      \item \textbf{Development of a New AI Scientist}: We developed Jr. AI Scientist, a new system that starts from a baseline paper and its associated codebase, and is capable of handling complex, multi-file implementations, overcoming a major limitation of previous AI Scientist systems.
      \item \textbf{Revealing Strengths and Limitations of Jr. AI Scientist}: We conducted extensive evaluations using open-source AI reviewers, Agents4Science, and author evaluation. The results demonstrate that Jr. AI Scientist generates higher-quality research papers than existing AI Scientists, while also revealing key challenges for future improvement.
     \item \textbf{Thorough Risk Report}: We report the observed risks throughout the project. We believe these reports offer insights into the areas that still require human expertise and the risks that may emerge as these systems evolve.
\end{itemize}















